\section{Infinite component groups and boundary attainment}\label{sec:profinite55}
This section retains independent finite-dimensional proofs of the strict-radius and exact translate-rank conclusions. Its output spaces are finite-dimensional. Let $H$ again be an arbitrary compact Hausdorff group, and set $K=H^\circ$. The quotient $H/K$ is profinite. We use two standard compact-group consequences of this fact and of Peter--Weyl theory \cite{Folland}: open normal subgroups form a neighborhood basis in a profinite group, and, for a continuous surjection $\pi:H\to G$ onto a compact Lie group, one has $\pi(K)=G^\circ$. To see the latter assertion, $\pi(K)$ is connected and normal, while $G/\pi(K)$ is a continuous group quotient of the profinite group $H/K$ and is therefore totally disconnected. This forces $G^\circ\subset\pi(K)$; the reverse inclusion is immediate. A continuous finite-dimensional representation of a profinite group has finite image, by the closed subgroup theorem and the no-small-subgroups property of a Lie group \cite{Hall}: an open normal subgroup can be placed inside a no-small-subgroups neighborhood, so its image is trivial.

Define the component radius
\begin{equation}\label{eq:R055}
                 R_0=\max_{s,j,g\in H}\rad_{Y_j}(F_{s,j}(gK)).
\end{equation}
Unlike $r(U)$, this quantity need not itself be realized by any finite quotient.

\begin{lemma}[Approximation by finite quotient radii]\label{lem:infradius55}
With $r(U)$ as in \eqref{eq:quotientradius55},
\[
                 \inf_{U\text{ open normal in }H}r(U)=R_0.
\]
\end{lemma}
\begin{proof}
Every open subgroup contains $K$, so $r(U)\ge R_0$. Fix $\eta>0$. Uniform continuity, for the finite family of interfaces, gives an identity neighborhood $V$ such that
\[
 \norm{F_{s,j}(xv)-F_{s,j}(x)}_j<\eta
       \quad(x\in H,\ v\in V,\ s,j).
\]
There is an open normal subgroup $U$ containing $K$ and contained in $KV$, by the profinite quotient basis. Write each $u\in U$ as $kv$ with $k\in K$, $v\in V$. Then $F_{s,j}(gu)$ lies within $\eta$ of $F_{s,j}(gk)$. A center for the latter component image consequently covers the former coset image with an additional error at most $\eta$. Thus $r(U)\le R_0+\eta$.
\end{proof}

\begin{lemma}[A legal finite-dimensional interface approximation]\label{lem:Lieapprox55}
For every $\eta>0$ there is a continuous homomorphism $\pi:H\to G\subset O(D)$ onto a compact Lie group and continuous maps $\widetilde F_{s,j}:G\to Y_j$ such that
\[
           \max_{s,j}\norm{F_{s,j}-\widetilde F_{s,j}\circ\pi}_\infty<\eta.
\]
Moreover, for every $g\in H$,
\[
 \left|\rad_{Y_j}(F_{s,j}(gK))-
 \rad_{Y_j}(\widetilde F_{s,j}(\pi(g)G^\circ))\right|<\eta.
\]
\end{lemma}
\begin{proof}
Approximate all coordinates of all $F_{s,j}$ by finitely many real representative functions. The direct sum of their orthogonal representations is $\pi$, and the approximants are continuous functions on its compact image $G$. Choose Euclidean inner products on the finitely many output spaces. Project each approximant to the closed convex set $Y_j$ by nearest-point projection. This projection is continuous and nonexpansive in the chosen Euclidean norm. Choose constants $a_j,b_j>0$ with $a_j\|v\|_j\le|v|_2\le b_j\|v\|_j$. An initial error less than $a_j\eta/b_j$ in $\|\cdot\|_j$ then gives a projected error less than $\eta$, by Euclidean nonexpansiveness. These constants are fixed before choosing the approximants. The closed subgroup $G\subset O(D)$ is a compact Lie group. Since $\pi(K)=G^\circ$, the two indicated images have Hausdorff distance less than $\eta$. Their constrained radii therefore differ by less than $\eta$.
\end{proof}

\begin{theorem}[Arbitrary compact groups on both strict sides]\label{thm:profinite55}
For any compact-group interface, $\varepsilon>R_0$ implies a finite stationary permutation realization. At $\varepsilon=R_0$, such a realization exists exactly when $r(U)\le R_0$ for some open normal $U$.

Suppose in addition that executable exact returns have all sufficiently large lengths. If $\varepsilon<R_0$, then for each $k$ there is a horizon-independent upper bound on the number of cuts of width at most $k$, with all intervening widths unrestricted. In particular $W_{N,\varepsilon}\to\infty$. No finite-component assumption is made in either assertion.
\end{theorem}
\begin{proof}
The stationary assertions follow from Lemma~\ref{lem:infradius55} and Theorem~\ref{thm:finitequotient55}. For the lower bound choose a component image of radius greater than $\varepsilon$ and then $\eta>0$ with $2\eta$ smaller than that strict difference. Apply Lemma~\ref{lem:Lieapprox55}. A machine correct for the original interface at error $\varepsilon$ is correct for the approximating interface at error at most $\varepsilon+\eta$. The chosen component radius of that interface is greater than $\varepsilon+\eta$. The projected commands still have cofinite exact returns and generate $G$ densely. Theorem~\ref{thm:radius54} therefore gives the desired occupation bound for the very same registers. This argument never assumes that $K$ is open or that executable products lying exactly in $K$ are dense.
\end{proof}
Theorem~\ref{thm:completeboundary56} now identifies this stationary boundary criterion with clocked boundedness, including $R_0>0$, and under the weaker approximate-return hypothesis. The zero-error boundary also admits the following independent translate-rank argument.

\begin{lemma}[Evaluation rank and translate dimension]\label{lem:evaluation56}
For a real function $f$ on a group, the supremum of the ranks of finite matrices $(f(y_lx_i))_{i,l}$ equals the dimension, possibly infinite, of the span of the functions $x\mapsto f(yx)$. These are left translates with the convention $(L_gf)(x)=f(g^{-1}x)$, after replacing $y$ by $g^{-1}$.
\end{lemma}
\begin{proof}
A finite-dimensional span bounds every evaluation rank. Conversely, if $f_1,\ldots,f_m$ are independent functions, their evaluation vectors $(f_1(x),\ldots,f_m(x))$ span $\R^m$: a proper span would have a nonzero annihilator, giving a linear dependence of the functions. Select $m$ evaluation vectors forming a basis. This gives an invertible evaluation matrix and proves the reverse inequality for every finite $m$.
\end{proof}

\begin{theorem}[The exact clocked boundary]\label{thm:zero55}
Assume cofinite executable exact returns. The following are equivalent:
\begin{enumerate}
\item $\sup_NW_{N,0}<\infty$;
\item all $F_{s,j}$ factor through one finite continuous quotient $H/U$, with $U$ open normal;
\item there is an exact stationary finite permutation realization.
\end{enumerate}
This includes groups with infinitely many connected components.
\end{theorem}
\begin{proof}
The implications from the second assertion to the third and from the third to the first are immediate. Suppose the first holds with bound $k$. Theorem~\ref{thm:profinite55} implies $R_0=0$, so every interface is constant on each $K$-coset.

Fix a real coordinate $f$ of one $F_{s,j}$. For any finite families of executable prefix products $x_i$ and suffix products $y_l$, use cofinite returns to pad all selected prefix words to one common length $A$ and all suffix words to one common length $B$. Choose an exact realization at horizon $A+B$ of width at most $k$. At its cut $A$ the matrix
\[
                            (f(y_lx_i))_{i,l}
\]
factors through the $k$ hidden labels: the left factors are the actual state distributions after the padded prefixes and the right factors are the actual conditional terminal expectations on the padded suffixes. Hence this matrix has rank at most $k$. Padding does not need to preserve hidden states; it only preserves the physical products used in the displayed entries. By density and continuity the same rank bound holds for arbitrary $x_i,y_l\in H$.

It follows that the span of the left translates of $f$ in $C(H,\R)$ has dimension at most $k$. Otherwise $k+1$ linearly independent translates would have an invertible evaluation matrix at some $k+1$ points, contrary to the preceding rank bound. Lemma~\ref{lem:evaluation56} formalizes this evaluation assertion. Equivalently it follows inductively: a nonzero function is nonzero somewhere, and a function outside the span of the preceding ones cannot have all its evaluations forced by theirs.

Take the sum $V$ of these translate spaces over all seed--query coordinates. It is finite-dimensional and invariant under left translation. Translation is continuous in the uniform norm and preserves that norm, so it gives a continuous representation of $H$ on $V$ with compact image; Haar averaging makes it orthogonal. Because $K$ is normal and every original coordinate is constant on $K$-cosets, $K$ acts trivially on every translate and thus on $V$. The representation factors through $H/K$. Its image is finite by the profinite finite-dimensional representation fact above. Its kernel $U$ is open normal. Since every original coordinate belongs to $V$, all interfaces are $U$-invariant and descend to $H/U$.
\end{proof}

\begin{proposition}[A nonattained zero-radius boundary]\label{prop:adic55}
On a compact group with infinitely many components, every positive error can have a bounded stationary realization while the zero-error clocked width is unbounded.
\end{proposition}
\begin{proof}
Let $H=\Z_2$ be the additive group of $2$-adic integers, with commands $0$ and $1$. Their positive products are dense, and the zero command supplies exact returns of every length. If $x=\sum_{n\ge0}x_n2^n$, $x_n\in\{0,1\}$, prescribe the binary success probability
\[
                         f(x)=\sum_{n\ge0}\frac{2x_n}{3^{n+1}}.
\]
Uniform convergence makes $f$ continuous. It is injective: at the first differing digit the contribution $2/3^{n+1}$ exceeds the sum $1/3^{n+1}$ of all possible later differences. Since $H$ is totally disconnected, $K=\{0\}$ and $R_0=0$. Truncation after $m$ digits factors through $\Z/2^m\Z$ and has uniform error at most $\sum_{n=m}^\infty2/3^{n+1}=3^{-m}$. Thus every positive tolerance admits a finite quotient machine. But the injective function $f$ cannot factor through any finite quotient. Theorem~\ref{thm:zero55} excludes bounded exact clocked width, and Theorem~\ref{thm:finitequotient55} excludes a finite exact stationary realization.
\end{proof}

\begin{theorem}[Linear exact width at a profinite boundary]\label{thm:adiclinear55}
For the interface of Proposition~\ref{prop:adic55},
\[
                         W_{N,0}=\Theta(N).
\]
More precisely, if $n$ is a power of two with $2n-1\le N$, then $W_{N,0}\ge n$, while $W_{N,0}\le N+1$. Every exact realization, with unrestricted intermediate widths, has at most $4k$ cuts in $\{1,\ldots,N\}$ of width at most $k$. Nevertheless, for each $\varepsilon>0$, $\sup_NW_{N,\varepsilon}<\infty$.
\end{theorem}
\begin{proof}
For nonnegative integers $z$ write $v_2(z+1)$ for the exponent of two in $z+1$. Incrementing the binary expansion, with $v_2(z+1)$ trailing ones, gives
\begin{equation}\label{eq:adicdifference55}
 b(z):=f(z+1)-f(z)=\frac53\,3^{-v_2(z+1)}-1.
\end{equation}
Let $n=2^m$, and form $B_n=(b(i+j))_{0\le i,j<n}$. Define a function on $\Z/n\Z$ by
\[
 \psi_m(x)=3^{-v_2(x)}\ (x\ne0),\qquad \psi_m(0)=3^{-m}.
\]
Here the valuation of a nonzero residue is independent of its representative. Since $1\le i+j+1\le2n-1$, its valuation equals that used in $\psi_m(i+j+1\bmod n)$, including the case $i+j+1=n$. Thus replacing row $i$ by the old row $(-i-1)\bmod n$ changes its entry at column $j$ to the kernel value at $j-i\bmod n$. This is the circulant with kernel $(5/3)\psi_m-1$.

On $\Z/n\Z$ one has the elementary identity
\[
            \psi_m=1-2\sum_{\ell=1}^m3^{-\ell}
                               \mathbf1_{2^\ell\mid x}.
\]
Use the eigenvector with entries $\exp(2\pi \mathrm i qj/n)$; its eigenvalue is the kernel sum against $\exp(2\pi \mathrm i qx/n)$. The sign is immaterial for subgroup indicators, whose transforms below are real. The Fourier eigenvalue at frequency zero is
\[
 \mu_0=\frac23 n6^{-m}>0.
\]
At frequency $q\ne0$ it is
\[
 \mu_q=-\frac{10n}{3}
           \sum_{\substack{1\le\ell\le m\\ n/2^\ell\mid q}}6^{-\ell}<0.
\]
Indeed a subgroup indicator has Fourier sum $n/2^\ell$ exactly when $n/2^\ell$ divides $q$, and zero otherwise. For $q\ne0$ the summand $\ell=m$ is always present. At $m=0$ the sole eigenvalue is $2/3$. Consequently $B_n$ is invertible for every dyadic $n$.

At an arbitrary command cut $t<N$, take a dyadic $n\le\min(t+1,N-t)$. For each $i<n$ choose a prefix of length $t$ with $i$ increment commands. For each $j<n$ choose two suffixes of length $N-t$ with respectively $j+1$ and $j$ increment commands. Fill all remaining positions with zero commands. The difference of their exact outputs is $b(i+j)$. Hence $B_n$ factors through the actual register $S_t$: one factor consists of the prefix state rows and the other of the differences of the two conditional suffix decoder columns. Invertibility gives $|S_t|\ge n$. The difference column need not be a legal decoder; it is used only for this rank bound on two legal experiments.

Taking $t=n-1$ proves $W_{N,0}\ge n$ whenever $2n-1\le N$. The largest dyadic $n\le(N+1)/2$ is greater than $(N+1)/4$, proving the linear lower order. Storing the number of increment commands uses at most $N+1$ labels and gives the exact upper bound.

For occupation, if $t<N$ and $|S_t|\le k$, the largest dyadic integer not exceeding $\min(t+1,N-t)$ is at most $k$, so $\min(t+1,N-t)<2k$. Such cuts lie within fewer than $2k$ positions of one of the two ends. Including the terminal cut gives at most $4k$ positive cuts. No width at another cut was restricted. The positive-error assertion is the finite-quotient truncation in Proposition~\ref{prop:adic55}.
\end{proof}
This supplies a sharp quantitative law on a genuinely infinite profinite group. Its linear exact-width obstruction and bounded positive-error realizations concern one fixed continuous interface; they are not a uniform approximation statement over all continuous outputs on that group.
